\section{Por qué los LLM son tan emocionantes}

\subsection{El crecimiento explosivo del tamaño de los parámetros}

Si los modelos de lenguaje n-gram ya existían desde la década de 1980, ¿por qué la gente no se entusiasmó con los modelos de lenguaje sino hasta años recientes? Uno de los cambios clave es el crecimiento exponencial del \textbf{tamaño de los parámetros}:

\begin{table}[H]
\centering
\begin{tabular}{lll}
\toprule
\textbf{Modelo} & \textbf{Año} & \textbf{Número de parámetros} \\
\midrule
BERT Large     & 2018 & 340 millones (340M)  \\
GPT-3          & 2020 & 175 mil millones (175B) \\
Modelos de vanguardia contemporáneos   & 2024--2025 & del orden de los billones (trillions) \\
\bottomrule
\end{tabular}
\caption{Tendencia de crecimiento del tamaño de los parámetros en los modelos de lenguaje}
\end{table}

Los parámetros pueden entenderse como los pesos de una red neuronal, en analogía con las conexiones sinápticas del cerebro. Más parámetros significan que el modelo tiene una mayor capacidad para «memorizar» y «comprender» información sobre el mundo.

\subsection{La expansión de la ventana de contexto}

Otro cambio clave es el crecimiento drástico de la \textbf{ventana de contexto} (context window):

\begin{table}[H]
\centering
\begin{tabular}{lll}
\toprule
\textbf{Arquitectura del modelo} & \textbf{Longitud típica de contexto} & \textbf{Orden de magnitud} \\
\midrule
Modelo n-gram       & $\sim$4 palabras          & unidades \\
RNN básica           & $\sim$20 palabras         & decenas \\
LSTM               & $\sim$200 palabras        & centenas \\
Transformer temprano   & $\sim$2,048 tokens    & millares \\
Gemini (2024+)    & $\sim$2,000,000 tokens & millones \\
\bottomrule
\end{tabular}
\caption{Evolución histórica de la ventana de contexto}
\end{table}

El problema del ciclo de probabilidades del modelo n-gram mencionado antes se debía precisamente a que su ventana de contexto tenía solo 4 palabras. Cuando la ventana se amplía al orden de millones de tokens, la cantidad de información que el modelo puede «ver» y procesar experimenta un cambio cualitativo.

\subsection{El surgimiento del Few-Shot Learning}

\begin{importantbox}{El hallazgo central del artículo de GPT-3}
El artículo de GPT-3 de 2020 («Language Models are Few-Shot Learners») reportó un comportamiento emergente clave (emergent behavior): cuando el número de parámetros del modelo alcanza alrededor de 175 mil millones, el modelo comienza a mostrar una capacidad exitosa de prompting zero-shot, one-shot y few-shot; sin ninguna actualización de gradientes (no gradient updates), y únicamente mediante las instrucciones del prompt o unos cuantos ejemplos, logra generalizar a tareas completamente nuevas.
\end{importantbox}

Las definiciones de estas tres formas de prompting:

\begin{itemize}
    \item \textbf{Zero-shot}: Solo se da la instrucción, sin ejemplos. Por ejemplo: «Translate English to French. Cheese:»
    \item \textbf{One-shot}: Se da un ejemplo. Por ejemplo: «Sea otter $\rightarrow$ Loutre de mer. Cheese $\rightarrow$»
    \item \textbf{Few-shot}: Se dan varios ejemplos (generalmente entre 3 y 5)
\end{itemize}

La gráfica central del artículo de GPT-3 muestra que, con un número bajo de parámetros, el desempeño zero/one/few-shot es cercano al azar; una vez que el número de parámetros supera cierto punto crítico, el desempeño aumenta de forma abrupta. Esto significa que los modelos de lenguaje a gran escala adquirieron una capacidad similar a la humana de «generalizar a partir de pocos ejemplos»: basta con dar unos cuantos ejemplos para generalizar a situaciones nuevas.

\begin{knowledgebox}{El significado de «No Gradient Updates»}
En los métodos tradicionales, para que un modelo domine una tarea nueva es necesario actualizar los pesos mediante fine-tuning (gradient updates). El artículo de GPT-3 subrayó especialmente «no gradient updates», porque el few-shot learning se realiza por completo en el momento de la inferencia (inference time), sin necesidad de modificar los parámetros del modelo. Esto redujo enormemente el costo de aplicar un LLM a una tarea nueva: basta con redactar bien el prompt.
\end{knowledgebox}

El docente también señaló la contribución del artículo de Chinchilla: mediante una estrategia de entrenamiento más eficiente (una asignación óptima entre datos y cómputo), es posible alcanzar un desempeño similar con menos parámetros. Esto trae consigo un menor consumo de energía, una inferencia más rápida y un menor uso de memoria. Pero los investigadores pronto destinaron los recursos ahorrados a continuar escalando (scale up), llevando el número de parámetros del orden de los cientos de miles de millones al orden de los billones.

\subsection{Resumen de la sección}

La razón fundamental por la que los LLM son tan emocionantes es la superposición de tres grandes cambios: (1) el número de parámetros creció del orden de los millones al orden de los billones; (2) la ventana de contexto creció de unidades a millones; (3) a una escala suficientemente grande surgió la capacidad de few-shot learning, capaz de generalizar sin necesidad de fine-tuning. Estos tres factores impulsaron juntos el cambio cualitativo de un «torpe autocompletado» a un «solucionador universal de tareas».

